% Nouveau EGA V (cote 25): archive cote 25, site page 131 (PDF page 132).
% English translation of the checked French annotated master.
% AI-translated by GPT-6 Astra (OpenAI, ChatGPT Pro mode).
\scanpage{25/131}
\mspage{\hand{34}}

by $\underline{O}_{X_o}(D_o)$ on D$_o$. In case b), One can assume D has no fixed
points, and it is more convenient to interpret the situation in terms of
morphisms into P (cf. 4.1.), so that G becomes the functor
\underline{Hom}$_S$(X,P), and \hand{$\mathcal{V}$}$_{w_o}$ will be the space \var{\xout{\illegible{1}}}{}{typewriter strike-out}{ega25_25_131-v0001} Hom$_{\underline{O}_{X_o}}$$(f_o^{\hand{*}}(\underline{\hand{\Omega}}^{1}_{P_o/k}), \underline{O}_{P_o})$.
In either case, there will be a natural homomorphism
\hand{$\mathcal{V}$}$_{u_o}\otimes_k V$ \hand{$\to$} \hand{$\mathcal{V}$}$_{w_o}\otimes_k V$ (where \hand{$\mathcal{V}$}$_{u_o}$ = dual of $u_o^{\hand{*}}($\hand{$\Omega$}$^{1}_{F_o/k})$ ) expressing the passage
from the deviation \hand{$\delta$} to the corresponding deviation \hand{$\delta$}' by applying h,
and the injectivity of this map \var{\xout{will result}}{results}{typewriter strike-out}{ega25_25_131-v0002} again from the injectivity
of \hand{$\mathcal{V}$}$_{u_o}$\hand{$\to$} \hand{$\mathcal{V}$}$_{w_o}$ , which in turn \var{\xout{\unsure{will come}}}{\unsure{comes}}{typewriter strike-out}{ega25_25_131-v0003} from the fact that $h_o$: $F_o$\hand{$\to$}$G_o$ is a
monomorphism.
\var{\margin{$/$}}{\placedmargin{$/$}}{marginal note placed}{ega25_25_131-v0004}

In practice, it hardly seems possible to write this last
part of the proof without referring to the small calculations of par.2\var{\over{5}{\unsure{6|5}}}{\unsure{6|5}}{writing over}{ega25_25_131-v0005}
(which there is no question of repeating here in a particular case). One
will point out that this gives rise to no vicious circles, since par.
2\var{\over{5}{\unsure{6|5}}}{\unsure{6|5}}{writing over}{ega25_25_131-v0006} \var{\xout{xxx}}{}{typewriter strike-out}{ega25_25_131-v0007} and the calculations \var{\xout{xxxxx}}{}{typewriter strike-out}{ega25_25_131-v0008} developed there depend only on the sorites
of differential calculus in par.16, and moreover \hand{4}.3. will no longer
be used in Chap. IV, except perhaps in the two n$^{o}$s that follow.

The interest of \hand{4}.3.\var{\inserted{a)}}{a)}{insertion}{ega25_25_131-v0009} is to show that under the conditions \var{\deleted{of a)}}{}{deletion}{ega25_25_131-v0010}\var{\inserted{\unsure{stated}}}{\unsure{stated}}{insertion}{ega25_25_131-v0011},
the parametrizing projective bundle can be interpreted as a subfunctor
of Div$_{X/S}$ , which permits an intrinsic interpretation of the notion
of a class (up to isomorphism of the parametrizing \var{\xout{scheme}}{}{typewriter strike-out}{ega25_25_131-v0012} projective
bundle) of linear systems on X/S, e\var{\xout{n terms of a}}{}{typewriter strike-out}{ega25_25_131-v0013} as being a subfunctor
$\overset{\hand{\vee}}{\text{P}}$ of Div$_{X/S}$ satisfying certain properties (\unsure{namely}, is
representable by a projective bundle, and the family of divisors defined
by the canonical injection of the latter into Div$_{X/S}$ is linear in the sense
of n$^{o}$\var{\over{4}{3}}{3}{writing over}{ega25_25_131-v0014}), which is \var{\xout{\unsure{precisely}}}{}{typewriter strike-out}{ega25_25_131-v0015}\var{\inserted[typed]{essentially}}{essentially}{insertion}{ega25_25_131-v0016} the classical point of view (where a
linear system of divisors is defined as a \uline{set} of divisors

\mspage{\hand{131}}
\tnote{25-131-1}{The verb translated as ``comes'' retains the source uncertainty between t and d in its stem before the canceled ra. The uncertainty in ``stated'' concerns the source's alternatives ``énoncées'' and ``énoncés''. The isolated e before ``as being'' remains from the partly canceled French phrase ``en termes d'un''; it has no separate translated word. The abbreviation n with a superscript o means number.}
